\subsubsection{Philosophy of micro test}
Micro test is designed to measure the performance of a chunk of code segment by comparing the time it consumed. Currently, we have designed the following micro tests, each targeting a specific computational pattern to evaluate language runtime and memory management efficiency.

To minimum the influence of environment and focus on the effect of garbage collector itself. Following points should be focused.

\begin{itemize}
    \item \textbf{Consistent implementation}: All micro test are implemented from scratch without any third-party dependencies, their code has totally same logic and algorithm. To make sure that the performance difference can only be caused by language and its garbage collector(if exists) itself.
    \item \textbf{Multi-platform}: Micro test must be replicated in different platform. To reduce the influence of platform features.
    \item \textbf{Automatic}: Micro test must be completed automatic, completing all steps by Python script. To minimum the influence of artificial mistake.
\end{itemize}


\subsubsection{Programs}
\begin{enumerate}
  \item \textbf{Fibonacci} — An iterative implementation of the Fibonacci sequence to test integer arithmetic throughput and loop execution efficiency.
  \item \textbf{Pi} — A Monte Carlo simulation to approximate the value of Pi, focusing on floating-point computation and pseudo-random number generation.
  \item \textbf{Matrix} — Dense matrix multiplication using nested loops to stress test memory allocation, cache usage, and raw arithmetic performance.
\end{enumerate}

Each of these micro benchmarks isolates a specific type of workload — iterative integer computation, probabilistic floating-point simulation, and linear algebra respectively — allowing for detailed performance comparisons across languages with and without garbage collection.


\subsubsection{Multi-platform}
The framework is designed to be executed on diverse hardware architectures to assess how performance characteristics, particularly those related to GC, might vary:
\begin{itemize}
    \item Personal computer (e.g., Macbook Pro) 
    \item Cloud-based server (e.g., Linux server, x86 architecture) 
    \item Embedded platform (e.g., ESP32, representing resource-constrained devices) 
\end{itemize}
This multi-platform testing strategy aims to provide a broader understanding of GC performance implications across different computational scales and constraints.


\subsubsection{Automatic Workflow}

\subsubsection{Hyperfine: A util to measure run time}
The framework relies on \texttt{hyperfine}, a command-line benchmarking tool chosen for its robustness and features suitable for this research. Key capabilities leveraged include:
It is our base our other script.

\begin{itemize}
    \item \textbf{Statistical Analysis}: Automatically computes mean, median, standard deviation, minimum, and maximum execution times over multiple runs.
    \item \textbf{Warmup Runs}: Allows specifying a number of initial runs (\texttt{--warmup N}) whose results are discarded, mitigating skewed measurements due to initial setup costs like JIT compilation or cache population.
    \item \textbf{Run Count}: Enables setting the number of measured runs (\texttt{--runs N}) for statistical significance.
    \item \textbf{Output Formats}: Exports results in various formats, including CSV for data analysis and Markdown for easy reporting.
    \item \textbf{Inter-process Measurement}: Accurately measures the wall-clock time of external commands.
\end{itemize}
A typical invocation within the framework looks like this:
\begin{lstlisting}[language=bash, caption={Example hyperfine usage in benchmark.sh }, style=shstyle]
hyperfine \
  --warmup 3 \
  --runs 10 \
  --export-csv "benchmark_TASK.csv" \
  --export-markdown "benchmark_TASK.md" \
  "./Rust_TASK <args>" \
  "./Go_TASK <args>" \
  "java Java_TASK <args>" # Example, actual execution might use run.sh
\end{lstlisting}

\subsubsection{Script}
\begin{itemize}
    \item \texttt{setup.sh}: Prepares the testing environment by installing necessary dependencies such as the Rust compiler, Go compiler, Java Development Kit (JDK), and the `hyperfine` benchmarking tool. It includes OS detection for compatibility with macOS and Linux.
    \item \texttt{prepare.sh}: Compiles all Go, Java, and Rust source files found within the task subdirectories, creating executable binaries or class files. Optimization flags are used for compiled languages (e.g., `rustc -O` or `rustc -C opt-level=3`).
    \item \texttt{benchmark.sh}: Executes the core benchmarking logic for a specified task directory. It utilizes the `hyperfine` tool to run the compiled programs, collect timing data, perform statistical analysis (mean, standard deviation), handle warmup runs, and export results to CSV and Markdown formats.
    \item \texttt{run.sh}: A helper script, typically invoked by `hyperfine` within `benchmark.sh`, responsible for executing a specific language's compiled program for a given task.
    \item \texttt{clean.sh}: Removes compiled artifacts (executables, `.class` files) to ensure a clean state before subsequent runs.
    \item \texttt{graph.py}: A Python script that reads the CSV output generated by `hyperfine` and uses `matplotlib` and `pandas` to create bar charts visualizing the mean execution times for each language within a benchmark task, saving them as PNG images.
\end{itemize}